\section{Mechanisation in Lean 4}
\label{sec:mech}

The following parts of this paper are formalised in the Lean package \code{lem-lean}
\citep{kreinzlemlean}: the graph model of LEM on which the vocabulary rests, the typed
crossing of Definition~\ref{def:typed} with the two maps of
Section~\ref{sec:three-kinds} as worked examples, the small deontic model of
Section~\ref{sec:deontic}, and the model of reasoning systems of
Section~\ref{sec:kernel} with its results. Section~\ref{sec:mech-limits} says what is
not formalised. The listings in this section and in the appendix are taken unchanged
from version 0.4.0 of the package (commit \code{c111aaf}).

\subsection{Layers, nodes and typed edges}

The package follows the minimal formalisation given in LEM: a typed multigraph. Nodes are
indexed by layer, predicates are local to a layer, and edges form a type of their own
with a map to their edge type. Duration is a function on the causal edges alone. The
central definitions are these (the complete file is in
Appendix~\ref{app:basic}):

\lstinputlisting[linerange={32-36,45-48,70-88}]{code/Basic.lean}

\noindent
Three things follow from the types alone.

\begin{enumerate}[label=(\arabic*)]
  \item \emph{No cross-layer application of a predicate.} If a predicate is declared at
        one layer, applying it to a node of another layer is a type error. The layer is
        part of the node's type, so this holds even in a model whose two layers have the
        same carrier.
  \item \emph{No duration on a non-causal edge.} \code{CausalDuration.deltaT} takes a
        causal edge, which is an edge together with a proof that its type is causal. For
        a constitutive, classificatory or normative edge that proof does not exist, and
        a term that asks for its $\Dt$ is ill-typed.
  \item \emph{Duration belongs to the edge, not to the pair of nodes.} Because edges are
        a type of their own, two causal edges between the same two nodes are two edges
        and may carry different durations.
\end{enumerate}

\noindent
The first two are covered by negative tests: terms that Lean must reject, with the
expected error recorded. The third is a theorem.

\subsection{Typed crossings}
\label{sec:mech-crossing}

An edge of the graph above carries a label for its type and nothing else. A typed
crossing in the sense of Definition~\ref{def:typed} has two components, and the package
states the second as propositions that have to be proved:

\lstinputlisting[linerange={21-66}]{code/Crossing.lean}

\noindent
The structure \code{Crossing} itself holds only the kind and the relation. What the
crossing carries is not stored in it: each preservation statement is a separate theorem
about the crossing. The declared crossing together with its proved statements is what
Definition~\ref{def:typed} calls a typed crossing. \code{transport} is the reuse of a
claim. It takes the preservation proof as an
argument, so a claim established at the source cannot be asserted at the target without
one. \code{kind\_not\_determined} holds by construction: the kind is a field that the
relation does not constrain. The theorem records that design decision. It is not a
discovery about maps.

The two maps of Section~\ref{sec:three-kinds} are stated as crossings with declared
kinds. For the cast, addition is carried unconditionally, subtraction under its side
condition and not without it, and the order is invariant. The listing shows the
definition and the two results on subtraction; the rest is in
Appendix~\ref{app:crossing}.

\lstinputlisting[linerange={80-84,95-110}]{code/Crossing.lean}

\noindent
The reading of a 64-bit vector as a number is treated in the same way: addition is
carried while nothing overflows, and not otherwise. Both maps are proved injective, and their declared kinds differ. This is the
formal counterpart of the observation that two maps of the same structure can be used
in different roles.

\subsection{Reasoning traces}

The model of Section~\ref{sec:kernel} and its results are formalised as follows. The
proofs are included to show their size. Proposition~\ref{prop:blind} is
\code{valid\_blind\_to\_mechanism}, Proposition~\ref{prop:char} is
\code{trace\_only\_test\_iff} and Corollary~\ref{cor:split} is
\code{no\_trace\_only\_test\_of\_split}. In this file the reading of \code{produce},
\code{valid} and \code{faithful} as a causal, a classificatory and a representational
edge is given in comments. The types do not enforce it. The remaining results of the
file are in Appendix~\ref{app:uniqueness}.

\lstinputlisting[linerange={29-40,45-49,74-101}]{code/Trace.lean}

\subsection{What is checked}

Table~\ref{tab:mech} lists the content of the package. It has no dependency on Mathlib
and contains no \code{sorry}. Lean's \code{\#print axioms} reports at most
\code{propext}, one of Lean's three standard axioms; the results on reasoning traces,
the deontic results and the locality results use no axioms. The toolchain is pinned to
Lean~4.35.0, and the build, including the negative tests, runs in continuous
integration.

\begin{table}[t]
\centering\small
\begin{tabularx}{\textwidth}{@{}P{3.1cm}L@{}}
\toprule
File & Content \\
\midrule
\code{Basic.lean} &
  Layers, layer-indexed nodes, layer states, local predicates and relations, declared
  co-reference, the six edge types, the event graph, causal duration, the three deontic
  types, the type-level stratum. \\
\code{Duration.lean} &
  No duration for constitutive, classificatory or normative edges; the total,
  option-valued reading of $\Dt$ is undefined exactly on the non-causal edges; two
  parallel causal edges may carry different durations. \\
\code{Locality.lean} &
  A layer-respecting state assignment on the whole graph is the same thing as a family
  of layer states, and a predicate on the whole graph the same thing as a family of
  local predicates. \\
\code{Deontic.lean} &
  With three norm types the type alone settles the standing of an omission; with two it
  does not; the reduction of prohibition to obligation leaves every verdict on an
  omission unchanged (Section~\ref{sec:deontic}; statements in
  Appendix~\ref{app:deontic}). \\
\code{Crossing.lean} &
  A crossing with a declared kind; preservation, reflection and invariance of
  predicates; invariance of a relation; an operation carried under a side condition;
  reuse of a claim by a preservation proof. Two worked crossings, $\N \to \Z$ and
  \code{BitVec 64} $\to \N$, with their side conditions and counterexamples. \\
\code{Trace.lean} &
  Propositions~\ref{prop:blind} and~\ref{prop:char} and Corollary~\ref{cor:split}; the
  uniqueness-based results and the two-route example of
  Appendix~\ref{app:uniqueness}. \\
\code{Rejections.lean} &
  Negative tests: a predicate of one layer applied to a node of another, also where
  both layers share a carrier; $\Dt$ asked of a normative edge. \\
\bottomrule
\end{tabularx}
\caption{Content of the \code{lem-lean} package, version 0.4.0.}
\label{tab:mech}
\end{table}

\subsection{What the mechanisation does and does not establish}
\label{sec:mech-limits}

\paragraph{It is partial.} Of the three cases, only the first is stated as a typed
crossing in Lean, and for two maps and a handful of operations. The package does not
reproduce what Mathlib's bundled morphisms and cast lemmas provide, and it does not
connect to them. The refinement case of Section~\ref{sec:refinement} is not formalised
at all. The reasoning-system model of Section~\ref{sec:kernel} is formalised on its own
terms and not as an instance of the crossing structure. The event graph and the
crossing structure are also not connected: an edge of the graph does not carry
preservation statements.

\paragraph{It checks consequences, not fit.} The checker establishes that the stated
results follow from the stated definitions, and that a declaration, once made, is used
consistently. It does not establish that the definitions fit any real system, or that a
declared type is the right one. The encoding does not decide the type of any edge: types
are declared by the modeller. The tests by which LEM proposes to decide them are not
formalised, and LEM records that they have not been validated between raters.

The same holds for the model of reasoning systems. That the results of
Section~\ref{sec:kernel} are machine-checked says nothing about whether
Definition~\ref{def:rs} describes a given model well. That is a question about the
crossing from the formal model to the system modelled, a representational one, and the
kernel does not test it. The paper's thesis applies to its own artefact: verifying the
proofs verifies the proofs.
